%% body_intro.tex
%% Shared scientific source for the research and TMLR entry points.
\section{Introduction}\label{sec:intro}

A contextual bandit must learn which action to take from the rewards of the
actions it has already tried. An upper confidence bound (UCB) policy balances
learning and reward by adding an uncertainty bonus to each predicted reward,
then choosing an action with a large resulting score. The bonus matters: a
policy that understates uncertainty can stop exploring too soon, while an
excessively large bonus can overwhelm differences in predicted rewards.

With a nonlinear predictor, retraining changes how the model represents past
observations. For example, a recommender may refit its reward predictor after
each interaction, then recompute the gradients of all previously recommended
items. Their recorded rewards stay the same, but the matrix built from their
gradients changes. A narrower width computed from that new matrix does not by
itself establish greater statistical certainty. This paper asks when such a
change of matrix preserves a confidence guarantee, and what that guarantee
costs in an implemented policy.

\paragraph{Why the old and new matrices need different justifications.}
The usual sequential confidence argument uses a feature for each observation
that was fixed before its reward was seen. We call the matrix that accumulates
those collection-time features the \emph{reference matrix}. Its stored features
stay fixed, even as the predictor continues to learn. The \emph{current matrix}
instead recomputes the gradients of the same past samples at today's model
parameters. We call this recomputation \emph{relinearization} or \emph{replay};
it does not collect new rewards.

The distinction is statistical, not merely computational. Once the model has
learned from past rewards, its recomputed features can depend on those rewards.
The usual confidence argument therefore does not automatically apply to the
current matrix. In notation, a sample collected at round $s$ contributes
$\bq_s(a_s)=\nabla_\theta\mu_{\btheta_s}(x_s,a_s)$ to the reference matrix,
but contributes $\nabla_\theta\mu_{\btheta_t}(x_s,a_s)$ when replayed at
round $t$. The reference matrix retains the predictable rank-one update used by the
standard self-normalized confidence argument; the replayed matrix need not.

Existing methods make different choices here. NeuralUCB and NeuralTS store
gradient features as observations arrive
\citep{zhou2020neural,zhang2021neural}. Neural tangent kernel (NTK) methods can keep
uncertainty features at initialization \citep{kassraie2022neural}, while
neural-linear methods use linear uncertainty in a learned last-layer
representation \citep{riquelme2018deep}. Fixed network linearizations and
post-hoc Laplace approximations offer related ways to measure predictive
uncertainty \citep{maddox2021fast,daxberger2021laplace}. We study the additional
obligations created by replaying the historical gradients at the current
parameter.

\paragraph{The central result: valid transport can still be unhelpful.}
\emph{Confidence transport} means using a justified comparison between the
reference and current matrices to convert a reference confidence band into a
score based on current curvature. A \emph{certificate} is a bound that makes
such a comparison valid. The \emph{width} is the geometric part of the
uncertainty bonus; a confidence radius scales it, and a separate term covers
model bias.

Direct reference scoring is already available. Transport is worth studying
because it exposes what a policy must establish if it chooses to score with
current curvature instead. It supplies no new observations or additional
statistical information. In fact, with the same history, center, radius and
bias, the certified transported width is at least the direct reference width.
An exact actionwise correction recovers that reference width
(\Cref{prop:reference-width-floor}). This comparison leaves the ordering of
adaptive regret open: different bonuses can change actions, which changes the
data on which subsequent decisions are based.

Our controlled study makes the cost concrete. The evaluated Hessian-transport
policy uses bounds on prediction Hessians to certify the change of matrix. It
has higher mean regret than every comparator in every condition. Its
recorded optimism checks pass, but its regret bounds are larger than a simple
deterministic cap. Thus a confidence construction can be valid under its
assumptions, conservative in its bound, and poor as a decision rule. These are
separate questions, and the paper examines each one.

\paragraph{What the analysis provides.}
The theory connects a statistical confidence band to the score the algorithm
actually uses. \Cref{fig:transport-interfaces} separates the reference matrix
$\bVbar_t$, the current matrix $\bV_t$, an applied operator $\Calg_t$ that may
approximate it, and the numerical width $\hat s_t(a)$ returned by a solver.
The matrices are symmetric positive definite (SPD), so their inverse quadratic
forms define widths. Each step needs its own bound, available before the
selected reward is observed.

\begin{figure}[t]
 \centering
 % Analytic schematic, manually derived from transport_theory.tex:
% asm:abstract-transport, eq:transported-score-new, prop:reference-width-floor.
% Native body-size TikZ; requires only the existing TikZ/PGF dependency.
% Arrows denote conditional certificate comparisons, not causal effects.
\begin{tikzpicture}[x=1cm,y=1cm,
  every node/.style={font=\normalsize,inner sep=0pt},
  cert/.style={draw=blue!55!black,line width=.65pt,->},
  object/.style={draw=black!28,fill=black!2,rounded corners=2pt,
    minimum width=3.6cm,minimum height=1.25cm,align=center}]
  \path[use as bounding box] (0,-6.8) rectangle (16.2,.3);
  \node[anchor=west] at (0,0)
    {\textbf{Before the reward: conditional certificate comparisons}};

  \node[object] (reference) at (1.8,-1.15)
    {Frozen confidence\\[5pt]$\bVbar_t,\ \bar s_t(a)$};
  \node[object] (current) at (6,-1.15)
    {Exact current metric\\[5pt]$\bV_t,\ s_t(a)$};
  \node[object] (operator) at (10.2,-1.15)
    {Applied SPD operator\\[5pt]$\Calg_t,\ s_{\Calg,t}(a)$};
  \node[object] (width) at (14.4,-1.15)
    {Numerical score width\\[5pt]$\hat s_t(a)$};
  \draw[cert] (reference.east) -- node[above=5pt]{A} (current.west);
  \draw[cert] (current.east) -- node[above=5pt]{B} (operator.west);
  \draw[cert] (operator.east) -- node[above=5pt]{C} (width.west);

  \node[anchor=north west,align=left] at (0,-2.2)
    {\textbf{A\quad Metric transport}\\[8pt]
     $e^{-\bar D_t}\bVbar_t\preceq\bV_t
       \preceq e^{\bar D_t}\bVbar_t$\\[7pt]
     Certified path-length envelope $\bar D_t$.};
  \node[anchor=north west,align=left] at (5.5,-2.2)
    {\textbf{B\quad Operator comparison}\\[8pt]
     $\kappa_{-,t}\bV_t\preceq\Calg_t
       \preceq\kappa_{+,t}\bV_t$\\[7pt]
     Two-sided SPD comparison.};
  \node[anchor=north west,align=left] at (11,-2.2)
    {\textbf{C\quad Solver certificate}\\[8pt]
     $s_{\Calg,t}(a)\le\hat s_t(a)$\quad every $a$\\[7pt]
     $\hat s_t(a_t)\le\alpha_t s_{\Calg,t}(a_t)$\\[7pt]
     Sharpness: played action only.};

  \draw[black!28,line width=.5pt] (0,-4.6) -- (16.2,-4.6);
  \node[anchor=west] at (0,-4.95)
    {Selection and sharpness use the \textbf{same realized width map} $\hat s_t(\cdot)$.};
  \node[anchor=west] at (0,-5.57)
    {\textbf{Direct reference comparison}\quad Same history, center, radius, and bias.};
  \node[anchor=west] at (0,-6.11)
    {$w_t^{\rm tr}(a)=e^{\bar D_t/2}\sqrt{\kappa_{+,t}}\,\hat s_t(a)
       \ge\bar s_t(a)$;\quad exact actionwise correction gives $c_t(a)s_t(a)=\bar s_t(a)$.};
  \node[anchor=west] at (0,-6.65)
    {These comparisons do not order the regret of policies with different adaptive histories.};
\end{tikzpicture}

 \caption{The steps between a statistical confidence band and an implemented
 score. A bound on changing the matrix, a bound on approximating it, and a
 bound on the solver's output serve different purposes. The formal conditions
 are in \Cref{asm:abstract-transport}. Every scored action needs an upper-valid
 numerical width; only the selected action needs the stated sharpness bound.
 Both conditions refer to the same widths used to choose that action.}
 \label{fig:transport-interfaces}
\end{figure}

The transported-UCB theorem combines these bounds and gives a regret guarantee
in terms of the reference matrix's accumulated information
(\Cref{thm:transported-ucb}). A two-sided matrix comparison is essential for
this guarantee: one direction preserves the upper confidence score, and the
other controls the width of the action actually chosen. We obtain such
comparisons from established Thompson--Finsler geometry
\citep{thompson1963certain,nussbaum1994finsler} and give computable certificates
using changes in model gradients (\Cref{lem:factor-spd-path,cor:hessian-path-new}).
The geometry and the UCB potential argument are established tools. The
contribution is to state and connect the distinct statistical and computational
obligations of scoring with replayed curvature.

For squared loss, we also construct the center of the confidence band. A ridge
correction based on collection-time features adjusts the current nonlinear
prediction. The confidence band around this \emph{corrected center} remains valid along any
predictable parameter path in the certified region, with supplied Taylor and
misspecification error bounds (\Cref{cor:corrected-center-new}). It does not
require the nonlinear optimizer to converge. The error bounds still have a
cost: persistent model mismatch can make regret guarantees linear or vacuous.
The finite-dimensional and near-linear consequences are given in
\Cref{cor:finite-d-rate-new,cor:global-near-linear-new}.

The scaled-tanh model permits a simpler comparison because its gradient
directions are fixed and only their weights change. Comparing those weights
gives a horizon-independent endpoint certificate and a tighter category-based
version (\Cref{cor:tanh-relative-weights}). This specialization does not cover
general rotating neural gradients.

\paragraph{What the experiments establish.}
The main study uses a 29-parameter scaled-tanh model with dense matrices and
direct Cholesky solves in float64 (\Cref{sec:transport-followup}). It compares
four policies in three prespecified conditions, with 50 paired seeds per
condition. All nine pointwise paired intervals favor the comparators over
Hessian transport. Transport inflation dominates Hessian's bound; the reference
policy's own-trajectory bounds also remain vacuous. All policies, including
the uncertified naive control, pass the recorded optimism diagnostics. These
floating-point checks neither certify numerical enclosures nor establish that
the correction was necessary in the observed runs.

The conditions change the reward task as well as curvature. The results
therefore describe this fixed teacher, model and shared prediction-based
optimizer selection; they do not isolate a causal effect of drift. Subsequent
analyses explain scores on recorded histories and assess tighter certificates
there. They do not run alternative adaptive policies. The original near-linear
audit is reported separately in \Cref{sec:transport-instantiation}; its bound
is dominated by the base component rather than transport inflation.
Approximate operators, conjugate-gradient performance and end-to-end
scalability are outside the empirical study.

\paragraph{Reading the paper.}
\Cref{sec:abstract-transport} starts with the exact-matrix case, then states the
general theorem. \Cref{sec:corrected-center} explains how the confidence center
is constructed. Readers interested primarily in the findings can continue to
\Cref{sec:transport-followup} and return to the certificate constructions when
needed. Proofs and the full experimental protocol are in the appendices.
\iflegacyextras
Earlier diagnostics are retained at their historical scope in
\Cref{app:legacy-evidence}.
\else
Earlier repository studies use a different, one-sided theorem and do not execute
the transported score. They are excluded from the evidence for this paper
(\Cref{app:legacy-evidence}).
\fi
